% AUTO-GENERATED by paper/code/gen_snippets.py -- do not edit.
\begin{lstlisting}[language=Python,style=fdarsinput]
import numpy as np, pandas as pd
from fdars import Fdata
from fdars.pace_fpca import irreg_fdata_from_lists, pace_fpca
growth = pd.read_csv(data_path("growth.csv"), index_col=0)
ARG = growth.index.values.astype(float)
X = growth.values.T.astype(np.float64)
pc = Fdata(X, argvals=ARG).to_pc(n_comp=3)
print("FPCA scores shape:", pc["scores"].shape)
print("singular values:", np.round(pc["singular_values"], 2).tolist())
# Sparsify: keep 8 random ages (of 31) per child, then fit PACE.
np.random.seed(42)
idx = [np.sort(np.random.choice(31, 8, replace=False)) for _ in X]
irreg = irreg_fdata_from_lists(
    [ARG[i].tolist() for i in idx],
    [X[k, i].tolist() for k, i in enumerate(idx)])
pace = pace_fpca(irreg, ncomp=2, bandwidth=0.5, sigma2=0.25,
                 work_grid=ARG.tolist())
r = [abs(np.corrcoef(pace["scores"][:, k], pc["scores"][:, k])[0, 1])
     for k in range(2)]
rmse = np.sqrt(np.mean((pace["fitted"] - X) ** 2))
print("PACE vs dense score correlation:", np.round(r, 3).tolist())
print("PACE trajectory RMSE vs full data (cm):", round(float(rmse), 2))
\end{lstlisting}
\begin{lstlisting}[style=fdarsoutput]
FPCA scores shape: (93, 3)
singular values: [227.54, 93.02, 43.73]
PACE vs dense score correlation: [0.983, 0.889]
PACE trajectory RMSE vs full data (cm): 2.72
\end{lstlisting}
